\begin{figure}[!htb] 
	\centering
		\resizebox{\textwidth}{!}{
			\subcaptionbox{Portion de \SI{30}{\second} de la trajectoire en force (affichage sous-échantillonné à \SI{100}{\hertz}, pour la première session)\label{subfig_unp_force_test_trajectory_a}}
			{
				\tikzsetnextfilename{subfig_unp_force_test_trajectory_a}
				\begin{tikzpicture}
				\centering
				%				
				\begin{axis}[%
					name=ax1,
					width=.6\linewidth,
					height=.25\linewidth,
					scale only axis,    
					xlabel={temps (\si{\second})},
					ylabel={Force (\si{\newton})},
					axis x line = bottom,
					axis y line = left, 
					xmin=70,xmax=100,
					ylabel near ticks,
					]
					%
					\addplot [
						color=Orient,
						smooth,
						]
					table [col sep=comma, x=time, y=force] {misc/unperturbed_force_signal.csv};
					%	
					%define coordinates of rectangle
					\coordinate (r1) at (axis cs:91,-5.8);
					\coordinate (r2) at (axis cs:93,7.5);
					\coordinate (r3) at (axis cs:91,7.5);
					\coordinate (r4) at (axis cs:93,-5.8);			
					\draw (r1) rectangle (r2);
					%			
				\end{axis}
				%
				\begin{axis}[%
					name=ax2,
					width=.35\linewidth,
					height=.35\linewidth,
					scaled ticks=false,   
					xlabel={temps (\si{\milli\second})},
					axis x line = bottom,
					axis y line = left, 
					xmin=91,xmax=93,
					%ymin=.025,ymax=0.25,
					at={($(ax1.south east)+(1cm,0)$)},
					]
					%
					\addplot [
						color=Orient,
						smooth,
						%line width=1.pt,
						]
					table [domain=91:93, col sep=comma, x=time, y=force] {misc/unperturbed_force_signal.csv};
					%					
				\end{axis}
				% draw dashed lines from rectangle in first axis to corners of second
				\draw [dashed] (r4) -- (ax2.south west);
				\draw [dashed] (r2) -- (ax2.north west);				
				
				\end{tikzpicture}
			}
		}	
	
		\resizebox{\textwidth}{!}{
			\subcaptionbox{Caractéristiques d'amplitude des cycles\label{subfig_unp_force_test_trajectory_b}}
			{
				\tikzsetnextfilename{subfig_unp_force_test_trajectory_b}
				\begin{tikzpicture}
				\centering
				%				
				\begin{axis}[%
					width=.95\linewidth,
					height=.25\linewidth,
					scale only axis,    
					xlabel={cycles},
					ylabel={Amplitude (\si{\newton})},
					ylabel near ticks,
					axis x line = bottom,
					axis y line = left, 
					]
					%
					\addplot [
						color=Orient,
						solid,
						mark=asterisk,
						]
					table [col sep=comma, x expr=\coordindex,  y=mag_var] {misc/unperturbed_force_signal_properties.csv};
					%	
					%define coordinates of rectangle
					\coordinate (r1) at (axis cs:42,9);
					\coordinate (r2) at (axis cs:67,14);			
					\draw[dashed,rounded corners] (r1) rectangle (r2);
					\node[below, align=center,fill=white,yshift=-2mm] at (axis cs:54.5, 14) {\SIrange{70}{100}{\second}};			
				\end{axis}
				%							
				\end{tikzpicture}
			}
		}
	\caption[Exemple d'une trajectoire en force non-perturbée]{Exemple d'une trajectoire en force de \SI{4.5}{\minute}, en l'absence de perturbations et de retours haptiques, pour un total de plus de 220 cycles. Lors de l'interaction avec le robot, le bras de l'utilisateur est soumis à son propre poids ainsi qu'à la masse inertielle du robot.}
	\label{fig_unp_force_test_trajectory}
\end{figure}